\textit{Data.} All data are simulated. Lorenz-63 is integrated with DOP853 (rtol and atol $10^{-12}$) from $(-8,7,27)+\mathcal{N}(0,I)$ after a 5-time-unit burn-in, then sampled for 10 time units at $\Delta t=0.01$ (1001 samples). Gaussian noise with standard deviation equal to a fraction of the per-component standard deviation of the clean signal is added at levels 0, 0.5, 1, 2, 5, 10\%. Each seed (0--19) gives a new initial condition and noise draw; every system sees the same 20 trajectories per level (6 levels $\times$ 5 systems $\times$ 20 seeds $=600$ main runs).

\textit{Tuning.} For each system and noise level, the smoothing parameter and the STLSQ threshold $\lambda$ are chosen by grid search on 5 separate trajectories (seeds 1000--1004): highest exact-support rate, then lowest median coefficient error, then the largest threshold. A tie in the tuner is therefore equality of the support rate and of the median error on the five trajectories; it does not require identical models. The grids are: $\lambda\in\{0.05,0.1,0.2,0.4,0.6\}$; SG $w\in\{11,21,41,81,161\}$; spline $\mu\in\{10^{-9},\dots,10^{-3}\}$ (7 values); TV $\alpha\in\{10^{-5},\dots,10^{-1}\}$; weak $w\in\{0.3,0.6,1,1.5,2.5\}$ time units. Tuning uses the ground-truth coefficients (oracle-assisted, equal for all systems), so absolute rates are optimistic. Selected values often lie on a grid edge (SG $w=11$ at five of six levels; $\lambda=0.6$ in 20 of the 30 cells of Table~\ref{tab:tuned}). In five cells (FD at 5\%; FD, SG, spline and TV at 10\%) no configuration recovered the support of any tuning trajectory, and $\lambda$ was selected by the median error alone. All systems share the same data, metric code and defaults; no baseline received other settings.

\textit{Revision of the tie rule.} The first version of the tuner resolved ties by grid order, i.e. by the \emph{smallest} threshold. Ties decided $\lambda$ in 19 of the 30 cells, and an audit showed that the conclusions below 5\% noise and for FD at 5\% depended on this arbitrary choice. The rule was replaced by the largest-threshold rule above (same tuning seeds and grids, no test data, smoothing parameters unchanged), and the main runs, the width sweep and the exponent ablation were run again. The new rule was adopted after the first test results were known, so it is not pre-registered; Table~\ref{tab:tuned} lists both sets of thresholds, Table~\ref{tab:thr} reports every system at every threshold with both choices marked, and the first-version runs stay in the registry as superseded.

\textit{What a tie is.} A diagnostic on the tuning seeds (\texttt{method/tie\_check.py}, one run per noise level) compared the models of the tied configurations. In 17 of the 19 tie cells they are identical on all five trajectories. In the other two the tied thresholds agree only in rate and median: for FD at 5\% (0.05 and 0.1), the cell on which the H2 verdict turns, the models differ on 3 of the 5 trajectories, and for spline at 10\% (0.4 and 0.6) on 1 of 5. In both, the selected threshold also has the lower mean tuning error, and a tie-break by mean error would select the same configuration in all 30 cells; the FD threshold at 5\% is nonetheless a choice between different models made on five trajectories.

\begin{table}[t]\centering\caption{Tuned threshold $\lambda$ (in parentheses the value under the first tie rule where it differs); last row: mean seconds per trial over the 120 main runs of the system.}\label{tab:tuned}
\setlength{\tabcolsep}{4pt}\begin{tabular}{lrrrrr}\toprule
Noise & FD & SG & Spline & TV & Weak \\\midrule
0\% & 0.6 (0.05) & 0.6 (0.05) & 0.6 (0.05) & 0.6 (0.05) & 0.6 (0.05) \\
0.5\% & 0.6 (0.1) & 0.6 (0.1) & 0.6 (0.1) & 0.6 (0.1) & 0.6 (0.05) \\
1\% & 0.6 (0.2) & 0.6 (0.2) & 0.6 (0.2) & 0.6 (0.2) & 0.6 (0.1) \\
2\% & 0.4 & 0.4 & 0.4 & 0.4 & 0.6 (0.2) \\
5\% & 0.1 (0.05) & 0.4 & 0.4 & 0.4 & 0.6 \\
10\% & 0.6 (0.4) & 0.2 & 0.6 (0.4) & 0.2 & 0.6 \\\midrule
s/trial & 0.10 & 0.10 & 0.12 & 1.46 & 0.10 \\
\bottomrule\end{tabular}
\end{table}

\textit{Hypotheses and tests} (wording of \texttt{proposal.md}, written before the runs; tests are Welch's $t$-test from \texttt{rh compare}, $p<0.05$, no multiple-comparison correction). H1: at every noise level $\ge1\%$ the weak form has a higher exact-support rate than every derivative-based baseline (higher mean and $p<0.05$). H2: SG, spline and TV each have a lower mean coefficient error than FD at noise $\ge1\%$; refuted if any smoother is not lower at some level. H3: at 0\% noise all five systems recover the support in all seeds. H4: the weak-form advantage depends on the test-function width, and too narrow a window loses the advantage (descriptive width sweep at 2\% noise). H5: evaluating the library on the smoothed state (SG, spline) instead of the raw noisy state improves support recovery at 2\% noise (descriptive).

\textit{Compute.} Apple M1 Pro CPU, 2 threads; the runtime per trial (Table~\ref{tab:tuned}) includes data generation. The registry holds 1540 current experiment runs (600 main, 600 threshold-grid, 200 width, 60 exponent, 80 library-state), a smoke test, six tie-diagnostic runs and 1100 superseded first-version runs; wall-clock was about 30 minutes for the first version and 20 for the revision.
